\subsection{Implantación progresiva del backend Laravel}
\label{sec:transicion-laravel}

La implantación progresiva de Laravel/PHP conserva los identificadores M1--M12, sus dueños de datos, los contratos públicos y sus versiones, los eventos JSON, la identidad Keycloak y las reglas sin conexión. Se separan los consumidores de integración de los trabajos internos y se prueban identidad, trazas y reversión antes de habilitar cada ola. El Anexo 4-P detalla las tecnologías que realizan estas capacidades; PostgreSQL/PostGIS, RabbitMQ, SQS y los clientes mantienen sus contratos durante el despliegue.


La implantación comienza por fijar contratos, dueños de datos y un extracto conciliado del WMS de 2013 en Talca.

\begin{itemize}
  \item Concepción inicia sus saldos con un conteo físico, porque hoy controla su stock por planilla.
  El esquema PostgreSQL de destino lo define el modelo de datos de la solución, sin suponer el motor del sistema antiguo.
  \item Las migraciones Laravel son aditivas y compatibles con los lectores de la ola previa.
  \item M1/M2/M5 se ensayan primero en un sitio piloto y después se habilitan los demás sitios, módulos y trabajadores.
  El WMS antiguo deja de escribir en cuanto se habilita cada ola y se retira al cerrar la marcha blanca de Talca, con sus datos ya migrados y conciliados.
  \item En coexistencia, cada operación tiene un único escritor autorizado y no hay doble escritura de stock, cobros o DTE.
  Una reversión de software vuelve a la versión compatible anterior del nuevo servicio, con el mismo estado y el mismo escritor. No reactiva el WMS antiguo como escritor de stock.
  \item La contingencia manual del Subdocumento 3 conserva el picking y la evidencia, pero no permite ejecutar a mano los 96 despachos de 05:30 a 07:00.
  Por eso la decisión de revertir se toma antes de esa ventana, con guías válidas emitidas por el ERP, conciliación previa y autorización del CLIENTE.
  \item Todo evento entre versiones usa JSON versionado.
  Cada ola se acepta con pruebas de 14 horas en terreno, 24 horas por sitio con dos relevos, guía previa al despacho, 96 salidas en la ventana crítica y reconciliación sin pedidos duplicados.
\end{itemize}
